\chapter{CPD标准模型：状态、观测、控制与可检验相态}
\label{chp:cpd_standard_model}

本附录把正文中分散的状态、观测、控制、约束、记忆与相态概念压缩成一组可比较的计算接口。所谓“标准模型”在这里不是声称所有智能系统具有同一内部结构，更不是把人脑、语言模型、强化学习系统和社会组织宣布为严格同构；它只给出一套最低限度的问题清单，使不同载体能够在相同任务、观测和干预条件下接受检验。我们将明确每个对象的空间、更新尺度与可观测读出，区分定义、模型假设、经验主张和AgentOS工程设计，并说明何时压缩会丢失决定任务成败的机制差异。CPD因此是HSF-HD一般状态动力学中的一个简化模型族，而不是凌驾于领域机制之上的宇宙终极方程。

\section{建模对象与极简接口：从普遍同构退回任务相关比较}
\label{sec:cpd_meta_framework}

我们先固定讨论范围。设一个候选智能系统在离散时刻 $k\in\mathbb{N}$ 的内部状态为 $x_k\in X$，其中 $X$ 可以包含连续向量、离散结构、事件记录和版本信息，但这些分量必须分别类型化；环境状态为 $s_k\in S$，观测为 $y_k\in Y$，行动为 $a_k\in A$。系统接收外部输入 $u_k\in U$ 与任务规范 $q_k\in Q$，按带参数和结构版本的更新核
\[
  x_{k+1}\sim P_{\theta_k,G_k}(\,\cdot\mid x_k,y_k,u_k,q_k),
  \qquad a_k\sim\pi_{\theta_k,G_k}(\,\cdot\mid x_k,q_k)
\]
演化。这里 $\theta_k$ 表示连续参数，$G_k$ 表示局部结构网络；二者随时间改变时，必须显式进入更新，而不能沿用固定模型的稳定性结论。全局分布表示可写为 $z_k\in Z\subseteq\mathbb{R}^d$，它通过任务相关身份、绑定映射与读出接口同 $G_k$ 联合，但不能替代图中对象身份、来源和权限关系。

在此基础上，我们把CPD接口写成
\[
  \mathfrak{C}_{\mathrm{CPD}}
  =(X,\mathcal{O},\mathcal{C},\mathcal{K},\mathcal{M},\mathcal{R}),
\]
而不再用五个字母宣称穷尽一切心智现象。$X$ 是类型化状态空间；$\mathcal{O}:S\times X\to\mathcal{P}(Y)$ 是观测核，描述环境与内部状态如何产生带噪记录；$\mathcal{C}$ 是受约束控制器族，生成候选行动或内部调度；$\mathcal{K}$ 是约束集合与不确定性记录，包括权限、资源、截止期、未解决冲突和模型误差；$\mathcal{M}$ 是按写入速度、所有权和保持期限分层的记忆接口；$\mathcal{R}$ 是任务读出与评价协议。各对象所用单位不同：时延用时间，存储用字节或记录数，计算用操作量或设备时间，实际能耗才用焦耳，概率不确定性与信息熵也不与这些量互换。

这个六元组首先是定义，不是经验定律。进一步假设系统具有有限观测、可修订状态、受约束行动和跨时保持机制，才形成一个具体模型。再进一步声称某类接口能解释人脑、LLM或群体协作，则属于需要数据支持的经验假说。最后，把 $h(t)/E/\Theta$、$\Psi$、SPC、TCE、CCM、Boundary、Offloading 与外部记忆Me映射到这些接口，只是AgentOS的一种工程设计：例如 $h(t)$ 可实现快速分布状态，$E$ 可保存任务相关误差，Boundary可执行能力检查，Me可承担带来源的外部记录；其他智能系统完全可以使用不同的机制满足同一接口契约。

跨载体比较必须限定任务集合 $\mathcal{T}$、允许干预集合 $\mathcal{D}$、读出度量 $d_R$ 与容差 $\varepsilon$。对两个系统 $A$ 和 $B$，我们只在存在编码、干预翻译和读出映射，使二者在这些查询下满足
\[
  \sup_{\tau\in\mathcal{T},\,d\in\mathcal{D}}
  d_R\!\left(R_A(\tau,d),R_B(\tau,d)\right)\leq\varepsilon
\]
时，称它们在指定任务上近似实现相同功能。该结论不推出内部结构唯一，更不推出意识、感受质或社会规范相同。AlphaGo的搜索、LLM的条件生成和组织的制度决策可以共享“提出—评价—提交”的抽象循环，却在时间尺度、可干预变量、责任主体和失败后果上显著不同。若忽略这些差异，所谓极简模型会把真正需要解释的机制删掉。

模型是否有用，应由分层证据决定。语法层检查类型、版本和单位；动力学层检查给定步长、时延和扰动下的有界性；任务层比较成功率、校准误差、恢复时间和严重失败；治理层检查授权、拒绝、撤回与责任记录。局部目标下降只说明所选代理量改善，不等于任务成功；预测有效也不证明系统整体同构。CPD的价值正在于提供一张可删减、可替换、可证伪的接口表：某一分量若在消融后不影响声明任务，就不应被宣布为智能本质；若某领域需要额外状态与机制，也必须扩展模型，而不能让“极简”成为拒绝证据的理由。

还应设置拒绝条件：若两个系统只能在训练分布内给出相似输出，却在结构编辑、延迟扰动、目标变化或权限撤销后出现不同响应，就不能继续宣称它们实现同一动力学。相反，即使内部材料和算法完全不同，只要在预先登记的干预集合上维持读出、约束与恢复性质，也可以获得有限范围的工程可替换性。这个判据把“跨载体”从修辞改成实验协议，并为后续各分量的逐项定义提供边界。
